\begin{table}[htbp]
\centering
\caption{Robustness of Triple-Difference Estimates}
\label{tab:robustness}
\small
\begin{tabular}{lccc}
\hline\hline
Specification & Coefficient & SE & \textit{p}-value \\
\hline
Baseline (Peak Stringency) & -1.580 & (1.231) & [0.205] \\
Binary Treatment (Median Split) & -0.371* & (0.197) & [0.066] \\
Cumulative Stringency (Mar--Jun) & -1.948 & (1.408) & [0.173] \\
Excl.~Never-Lockdown States & -1.140 & (1.053) & [0.285] \\
Alt.~Comparison: CPT Professional & -1.429 & (1.373) & [0.303] \\
Monthly Stringency (Post Only) & 0.009** & (0.003) & [0.012] \\
Randomization Inference (1000 perms) & -1.580 & (1.142) & [0.176] \\
Placebo (March 2019) & -0.116 & (0.439) & [0.793] \\
\hline\hline
\end{tabular}
\begin{minipage}{0.85\textwidth}
\vspace{4pt}
\footnotesize
\textit{Notes:} All specifications use the same three-way fixed effects (state $\times$ service, service $\times$ month, state $\times$ month) and cluster standard errors at the state level. Outcome is log total paid. Baseline uses peak April 2020 stringency (standardized 0--1). Binary treatment splits states at median stringency. ``Monthly Stringency'' uses time-varying OxCGRT stringency (0--100 scale), sample restricted to April 2020--December 2022 (OxCGRT coverage period); note the different treatment scale compared to baseline. For RI, SE column reports the standard deviation of the permutation distribution; the $p$-value is the share of permuted coefficients with absolute value exceeding the actual. Placebo assigns treatment in March 2019 using pre-period data only.
\end{minipage}
\end{table}
